\section{Limitations}
\textbf{Benchmarks.} RadHAR has two subjects, no subject identifiers, and test
recordings from the same sessions as the training recordings, so it cannot
measure generalisation to new users. MiliPoint offers $xyz$ only and three test
subjects, and its action names are not published. \textbf{Baselines.} All
baselines were re-implemented by us under one training budget (at most 15/30
epochs with early stopping, CPU-only). Our PointNet++ did not train well on
MiliPoint's six-point frames, and DGCNN and Point Transformer were evaluated on
RadHAR only. Published implementations with tuned neighbourhood
parameters and longer training may perform better. \textbf{Robustness.}
Perturbations are applied to recorded point clouds and are not a physical radar
simulation. Under extreme sparsity (12.5\% of returns kept), STPT-Lite falls to
PointNet level on MiliPoint, and on RadHAR it is far more sensitive to injected
false detections than PointNet with the same robust training (55 vs.\ 96\% at
20\% outliers). \textbf{Deployment.} Latency was measured on one desktop CPU
thread, and storage is computed for INT8 weights. On-device deployment was not
performed. \textbf{Scope.} Single person, fixed 2-s windows, no fall labels,
and no clinical validation.

\section{Conclusion and Future Work}
We presented STPT-Lite, a factorised spatio-temporal point transformer for
privacy-preserving activity recognition from sparse mmWave radar point clouds,
together with Doppler-based person centring and clutter gating. On 49 MiliPoint
actions performed by unseen people, the distilled, noise-robust model reaches
\mistptlitekdrobustAcc\% top-1 (\mistptlitekdrobustTopthree\% top-3), about 17 points above the best re-implemented baseline, at
\oursParamsK\,K parameters, \oursMflops\,MFLOPs and \oursLat\,ms per window. On
RadHAR all models saturate at about 95--98\%, and the pre-processing, not the
architecture, explains the gain over the original 90.5\%. A dedicated audit
shows that random window splits leak most test frames into training, which
accounts for much of the spread of published MiliPoint results. Future work:
leave-one-subject-out and cross-room evaluation on datasets with subject
identifiers and fall events, deployment on the radar's companion processor, and
self-supervised pre-training on the large unlabelled portions of both datasets.
